% -----------------------------------------------------------------------------------------------------------------------------------------
% 第8章 Verdictによる実証：公共的価値接続とフィールド実験
% -----------------------------------------------------------------------------------------------------------------------------------------

\chapter{Verdictによる実証：公共的価値接続とフィールド実験}
\label{chap:integrated_governance}

% -----------------------------------------------------------------------------------------------------------------------------------------
\section{はじめに}
\label{sec:ch8_integrated_intro}

第7章では，秘密保護と民主的正当性を同時に満たす政策評価システムとして，市民討議による原則策定，Constitutional AIによる価値の埋め込み，LLM as a Judgeによる評価，隔離実行環境による秘密保護を組み合わせる設計を提示した．
しかし，この設計は一つの実証実験だけで検証できるものではない．
なぜなら，このシステムが応答しようとしている問題は，技術的な実行可能性だけでなく，公共的価値を機械可読な評価基準へ変換できるか，市民がその変換過程を正当なものとして受け入れるか，さらに秘密情報を扱う政策評価に適用できるかという複数の水準にまたがるからである．

本章では，政策評価システムVerdictを対象として，第7章の設計を三つの水準から検証する．
第一に，盛岡市総合計画市民アンケートを用いたラボ実証により，市民由来の公共的価値をLLM judgeの評価行動へ接続する方法を比較する．
第二に，岩手県内の地域計画・自治会計画を対象とするフィールド実験計画により，市民討議から評価原則を生成し，人間の採点を基準としてJudge出力を検証する手続きを設計する．
第三に，秘密保護については，当日のフィールド実験で全面的に検証するのではなく，教材用の秘匿ケースとラボ実験を組み合わせる限定的な検証として位置づける．

したがって本章の立場は，「Verdictが完成済みの行政システムとして現場投入された」と主張することではない．
むしろ，公共的価値を評価AIに接続する過程を，ラボ実験，フィールド実験，秘密保護実験に分解し，どの部分が検証済みで，どの部分が今後の課題として残るかを明確にすることにある．

% -----------------------------------------------------------------------------------------------------------------------------------------
\section{Verdictの検証対象}
\label{sec:ch8_verdict_scope}

\subsection{Debate--Train--Judgeの三段階}

Verdictは，Debate，Train，Judgeの三段階からなる政策評価プロセスである．
Debateは，市民や地域関係者が政策判断の原則を言語化する段階である．
Trainは，その原則を機械可読な評価基準へ変換する段階であり，プロンプト憲法，検索用原則レコード，LoRAアダプタなど複数の実装形態を取りうる．
Judgeは，政策案やケース文に対して，生成された基準に基づく評価を行う段階である．

第7章ではこの三段階を，秘密情報を含む政策提案の評価まで含むフルアーキテクチャとして提示した．
これに対し，本章の実証は，フルアーキテクチャを一度に検証するものではなく，次のように分割する．

\begin{table}[htbp]
\centering
\caption{Verdict実証における検証範囲}
\label{tab:ch8_verdict_validation_layers}
\begin{tabular}{p{0.20\linewidth}p{0.34\linewidth}p{0.34\linewidth}}
\toprule
検証水準 & 検証する問い & 本章での扱い \\
\midrule
ラボ実証 & 市民由来の原則はLLM judgeの判断を変えるか & 盛岡市民アンケート由来の100ケースでC0--C3を比較する \\
フィールド実験 & 住民討議から評価原則を作り，人間採点を基準化できるか & 岩手県内の地域計画・自治会計画を対象に80分ワークショップを設計する \\
秘密保護 & 秘密情報を含む提案を評価しつつ漏洩を抑えられるか & 教材用秘匿ケースとラボ実験で限定的に検証する \\
\bottomrule
\end{tabular}
\end{table}

この分割により，第7章の二重制約を過大に実証したかのように見せることを避ける．
フィールド実験が主に検証するのは，民主的正当性，すなわち評価原則が地域の参加者から生成され，人間の採点と比較可能な形で記録されるかである．
秘密保護と閉域実行の検証は，別途の技術実験および教材用ケースによって補完される．

\subsection{LoRAをめぐる位置づけ}

VerdictのTrain段階では，市民由来の価値を評価器に接続する方法として複数の選択肢がある．
本研究で比較するのは，C0（地域固有の原則を入れないベースJudge），C1（市民由来の原則レコードを参照文脈として与える条件），C2（原則を凝縮したプロンプト憲法を与える条件），C3（ラベル付き事例でLoRA微調整を行う条件）である．

ここで重要なのは，LoRAを「常に最も高精度な方法」として位置づけないことである．
本章で示す実験では，小規模な訓練事例数のもとでLoRAが最良の即時精度を示すわけではない．
むしろLoRAの意義は，地域の価値プロファイルを小さなアダプタとしてパッケージ化し，自治体・政策領域ごとに版管理し，閉域環境で運用できる点にある．
したがって，本章では，評価精度としての効果と，秘密保護・運用管理としての効果を分けて論じる．

% -----------------------------------------------------------------------------------------------------------------------------------------
\section{ラボ実証：市民由来原則のJudge接続}
\label{sec:ch8_lab_benchmark}

\subsection{データと評価課題}

ラボ実証では，盛岡市総合計画市民アンケートを素材として，市民の選好や政策判断に関する記述を原則レコード，プロンプト憲法，LoRA訓練事例，評価用ケースに変換した．
アンケートは，公共交通，インフラ，福祉，参加，地域コミュニティ，環境，経済活力など複数の政策領域を含む．
評価課題は，二つの政策選択肢のうち，市民アンケートに現れた公共的価値により整合する選択肢を選ぶ二択判断として構成した．

評価セットは100件の保持ケースからなり，各ケースはアンケート設問または自由記述クラスタに紐づく根拠を持つ．
正解ラベルは研究者が作成したものであり，Aが33件，Bが67件である．
このラベル分布では単純な正解率が誤解を招くため，主指標にはBalanced Normative Fidelity（BalNF）を用いた．
BalNFはAラベルとBラベルそれぞれの正答率を平均する指標であり，全てBを選ぶだけのモデルは単純正解率0.67であってもBalNFは0.50にとどまる．

\subsection{比較条件}

比較条件は次の四つである．

\begin{description}
\item[C0 Base] ケース文のみを与え，盛岡固有の原則を与えない条件である．
\item[C1 Retrieval] 市民由来の原則レコードを参照文脈としてプロンプトに付与する条件である．本実装では検索ランキングを最適化したRAGではなく，原則文脈を制御して与える条件として扱う．
\item[C2 Prompt Constitution] 同じ原則レコードを凝縮し，評価ルーブリックとしてシステムプロンプトに注入する条件である．
\item[C3 LoRA] ラベル付き政策選択事例を用いてLoRAアダプタを訓練し，評価器の重みに地域価値を埋め込む条件である．
\end{description}

評価対象モデルは，Gemma3 instruction-tunedモデル（1B，4B，12B，27B）とQwen3 denseモデル（0.6B，1.7B，4B，32B）である．
C3では77件の訓練事例を用いた．
各条件の比較は，同一の保持ケース集合に対して行い，95\%ブートストラップ信頼区間を算出した．

\subsection{主要結果}

主要な結果を表\ref{tab:ch8_benchmark_results}に示す．

\begin{table}[htbp]
\centering
\caption{盛岡市民アンケート由来100ケースにおけるBalNF}
\label{tab:ch8_benchmark_results}
\begin{tabular}{lcccc}
\toprule
モデル & C0 & C1 & C2 & C3 \\
\midrule
Gemma3-1B  & 0.587 & 0.632 & 0.536 & 0.572 \\
Gemma3-4B  & 0.632 & 0.645 & \textbf{0.713} & 0.608 \\
Gemma3-12B & 0.729 & 0.653 & 0.675 & \textbf{0.740} \\
Gemma3-27B & \textbf{0.782} & 0.736 & 0.691 & 0.760 \\
Qwen3-0.6B & 0.572 & 0.564 & \textbf{0.586} & 0.550 \\
Qwen3-1.7B & 0.485 & 0.493 & \textbf{0.538} & 0.485 \\
Qwen3-4B   & 0.721 & 0.718 & \textbf{0.736} & 0.699 \\
Qwen3-32B  & 0.720 & \textbf{0.751} & 0.704 & 0.744 \\
\bottomrule
\end{tabular}
\end{table}

第一に，最も明確な要因はモデル規模である．
Gemma3系では，C0が1Bの0.587から27Bの0.782へ上昇しており，Qwen3系でも0.6Bから4Bにかけて大きく改善する．
これは，市民由来の原則をどのように接続するか以前に，評価器としての基礎能力が結果を大きく左右することを示している．

第二に，プロンプト憲法の効果は中規模モデルで最も見えやすい．
Gemma3-4BではC2が0.713に達し，C0の0.632から0.081ポイント改善した．
しかし，この効果は大規模モデルでは一貫して残らない．
Gemma3-27Bでは，C0が0.782と最も高く，C2は0.691に低下する．
この結果は，市民由来の原則が常に性能を上げるというより，中規模モデルにおいて明示的な評価ルーブリックが判断を支援する可能性を示すものと読むべきである．

第三に，77件の訓練事例によるLoRAは，最良の非訓練条件を一貫して上回らない．
Gemma3-12BではC3が0.740と高いが，Gemma3-27BやQwen3-4BではC0またはC2を下回る．
したがって，本実験から直ちに「LoRAが公共的価値接続の最良手法である」とは言えない．
一方で，追加のApple MLX感度分析では，根拠説明を含む事例で訓練したQwen3-4Bのq/vアダプタが，MLX上のC2ベースライン0.706から0.729へ方向的に改善した．
ただし，この差の信頼区間はゼロをまたぐため，決定的な効果ではなく，より良い訓練目的や事例数によって改善可能性が残るという限定的な示唆にとどまる．

\subsection{含意}

ラボ実証の含意は二つある．
第一に，市民由来の原則はLLM judgeの判断に影響を与えうるが，その効果はモデル規模，提示形式，訓練データ量に依存する．
第二に，C2とC3は同じ軸で比較すべきではない．
C2は即時の評価精度を得やすく，内容の確認や修正も容易である．
一方，C3は，地域の価値プロファイルをアダプタとして版管理し，自治体別・政策領域別に切り替え，閉域環境で運用するための制度的・運用的利点を持つ．
本研究にとって重要なのは，精度の勝者を一つ選ぶことではなく，公共的価値の接続方法を，評価性能，秘密保護，運用可能性の三軸で整理することである．

% -----------------------------------------------------------------------------------------------------------------------------------------
\section{フィールド実験：地域社会における原則生成}
\label{sec:ch8_field_experiment}

\subsection{対象と位置づけ}

フィールド実験は，岩手県内の自治体・自治会を対象として設計する．
ただし，盛岡都市圏の広域調整問題ではなく，盛岡都市圏を除く地域における農村計画，地域計画，総合戦略，自治会計画などを主な対象とする．
これは，地域の生活課題が明確であり，住民が評価原則を自分たちの言葉で形成しやすい領域を対象とするためである．

参加単位は自治会またはそれに準じる地域団体とし，参加者数は10〜20名程度を想定する．
実験はオンライン討議ではなく，対面の短時間ワークショップとして実施する．
この設計は，第7章で示したフルスケールの市民討議をそのまま実施するものではなく，現場負担を抑えつつ，Debate段階の中核である「住民由来の評価原則生成」と，人間採点による基準データ作成を検証する最小実装である．

\subsection{80分ワークショップの手順}

当日のワークショップは，前半40分，休憩10分，後半40分からなる．
前半では，地域の課題や将来像について参加者が発言し，ファシリテーターが発言を整理しながら，政策案を評価するための原則を3〜5本にまとめる．
ここで作成される原則は，VerdictにおけるDebateの出力であり，後続分析ではJSON形式の原則レコードとして扱う．

休憩時間には，ファシリテーターと地域役員が，前半で形成された原則をもとに二つのケース文案を作成する．
ケースは，地域が実際に直面しうる政策選択や事業案を簡潔に記述するものであり，参加者が後半で採点できる粒度に調整する．

後半では，参加者がケース文を確認し，各原則に照らして個人採点を行う．
採点後，短い共有を行い，どの原則が判断に効いたのか，どの点に迷いがあったのかを記録する．
この人間採点は，Judgeを代替するものではない．
むしろ，Judge出力の妥当性を検証するためのゴールドラベルであり，同時に参加者が評価プロセスを正当なものとして受け止めるかを確認する制度的データである．

\begin{table}[htbp]
\centering
\caption{フィールド実験プロトコル}
\label{tab:ch8_field_protocol}
\begin{tabular}{p{0.18\linewidth}p{0.30\linewidth}p{0.38\linewidth}}
\toprule
段階 & 内容 & 出力 \\
\midrule
前半40分 & 地域課題の共有と評価原則の作成 & 原則3〜5本，発言メモ \\
休憩10分 & ケース文案の作成 & 参加者由来のケース2件 \\
後半40分 & ケース確認，個人採点，短い共有 & 人間採点，迷い・理由の記録 \\
会後分析 & 原則・ケース・採点を用いたJudge評価 & Judge出力と人間採点の比較 \\
\bottomrule
\end{tabular}
\end{table}

\subsection{会後のTrainとJudge}

フィールド実験では，当日にAI評価を実行しない．
これは，参加者に対してAIによる即時判定を提示することが，討議そのものを誘導する可能性があるためである．
当日は，人間による原則形成と採点に集中し，会後に研究者が原則レコード，ケース文，人間採点を整理してTrainおよびJudgeを実行する．

会後分析では，原則レコードを用いてC1またはC2条件を構成し，必要に応じて小規模なC3アダプタを作成する．
Judgeの出力は，人間採点との一致，不一致，不一致が生じた理由の質的分析によって評価する．
特に，Judgeが多数派の採点に一致したかだけでなく，参加者が重視した原則を推論の中で保持しているか，少数意見や迷いの理由を消去していないかを確認する．

\subsection{評価指標}

フィールド実験の評価指標は，制度的指標，研究上の指標，技術的指標に分ける．
制度的指標としては，原則作成過程への満足度，採点手続きの理解度，AI評価に対する信頼と不安を測定する．
研究上の指標としては，人間採点とJudge出力の一致・不一致，原則ごとの判断理由，参加者が生成した原則の抽象度と具体性を分析する．
技術的指標としては，会後のTrain時間，Judge推論時間，出力形式の安定性，教材用秘匿ケースに対する秘密検出の成否を記録する．

参加地域には，原則リスト，ワークショップ記録，ケース別の論点整理，生成AI活用ガイドを返却する．
この返却物は，研究データ収集のためだけでなく，地域側が自分たちの計画づくりに使える成果物として設計する．

% -----------------------------------------------------------------------------------------------------------------------------------------
\section{秘密保護の限定検証}
\label{sec:ch8_confidentiality}

第7章で提示したVerdictの核心は，秘密情報を含む政策提案を評価しながら，その秘密を外部に漏らさない点にある．
しかし，実際の自治体・事業者の秘密情報をフィールド実験の初期段階で扱うことは，倫理的・制度的な負担が大きい．
そのため，本章のフィールド実験では，秘密保護の本格運用を直接検証するのではなく，教材用の架空ケースとラボ実験によって限定的に検証する．

教材用ケースでは，例えば事業者の採算情報，ルート案，未公開の交渉条件などを模した秘密フィールドを含む政策提案を用いる．
Judgeは，評価に必要な抽象的論点を出力しつつ，秘密フィールドそのものを出力に含めないことが求められる．
この出力に対して，秘密検出モジュールによる検査と，人間による確認を行う．

既に第7章で述べた予備実験では，Constitutional AI訓練によって漏洩率は低下したものの，ロールプレイ型の攻撃に対してはなお脆弱性が残った．
この結果は，LoRAやプロンプトによる論理的制御だけで秘密保護を達成できないことを示している．
したがって，Verdictの秘密保護は，AIの出力制御，秘密検出，隔離実行環境，評価結果のアクセス制御を組み合わせる多層防御として設計される必要がある．

この点は，生成AIを政策のメタネイチャーに接続する際の危うさとも関係する．
生成AIは，多数の知識・価値・事例を接続し，政策評価の条件を可視化する基盤になりうる．
しかし，基準の出所，訓練データ，出力の検証手続きが曖昧であれば，それは公共的判断を支える知識接続ではなく，もっともらしいスロップの大量生成に転化する．
Verdictの実証は，この危険を避けるために，市民由来の原則，人間採点，Judge出力，秘密検出を相互に照合可能な形で記録することを重視する．

% -----------------------------------------------------------------------------------------------------------------------------------------
\section{本章の限界}
\label{sec:ch8_limitations}

本章の実証には四つの限界がある．

第一に，ラボ実証の100ケースは，C0--C3の比較と信頼区間の算出には有用であるが，政策領域ごとの詳細な差異を分析するには十分ではない．
また，正解ラベルは研究者が作成しており，独立した複数ラベラーによる検証や一致度の測定が今後必要である．

第二に，C1およびC2の効果は，市民由来の原則そのものの効果と，単に構造化されたルーブリックを与えた効果を完全には分離できていない．
今後は，汎用的な市民参加ルーブリック，ランダムな原則，別地域の原則を用いた対照条件を追加する必要がある．

第三に，C3のLoRA実験は，77件の訓練事例という小規模条件での結果であり，LoRA一般の有効性を否定するものではない．
訓練事例数，学習率，rank，対象モジュール，根拠説明を含む訓練目的を変えた検証が必要である．

第四に，フィールド実験は，主に民主的正当性と評価基準生成の検証であり，秘密保護型マーケットサウンディングの実運用をそのまま検証するものではない．
秘密情報を含む実提案の評価，事業者参加，行政内部でのアクセス制御，法的位置づけは，次段階の実証課題として残る．

% -----------------------------------------------------------------------------------------------------------------------------------------
\section{小括}
\label{sec:ch8_summary}

本章では，第7章で提示したVerdictを，ラボ実証，フィールド実験，秘密保護の限定検証という三つの水準に分けて位置づけた．
盛岡市民アンケートを用いた100ケース実験では，市民由来の原則をLLM judgeへ接続する方法として，C0，C1，C2，C3を比較した．
結果として，モデル規模が最も大きな要因であり，プロンプト憲法は中規模モデルで有効に働く可能性がある一方，77件の訓練事例によるLoRAは最良の非訓練条件を一貫して上回らなかった．
ただし，LoRAは評価精度だけでなく，地域価値を閉域環境で版管理・運用するためのアダプタとして制度的意義を持つ．

フィールド実験については，岩手県内の地域計画・自治会計画を対象とする80分ワークショップを設計した．
当日は，住民が評価原則を作成し，ケースに対して人間採点を行う．
TrainとJudgeは会後に実行し，人間採点との比較によってJudge出力の妥当性を検証する．
この設計により，Verdictの実証は，単なるAI評価システムの性能評価ではなく，公共的価値がどのように生成され，記録され，機械評価に接続されるかを検証するものとなる．

以上より，本章は，第7章の設計論を，実証可能な検証単位へ分解した．
第9章では，この結果を踏まえ，本研究全体が自治体の政策形成能力，政策評価における知識接続，および政策のメタネイチャーの議論に対して持つ含意を整理する．

% -----------------------------------------------------------------------------------------------------------------------------------------
\printbibliography[heading=subbibliography,title={章末参考文献}]
